\documentclass[UTF8,12pt]{ctexart}

\usepackage[a4paper,margin=2.2cm]{geometry}
\usepackage{amsmath,amssymb}
\usepackage{booktabs}
\usepackage{enumitem}
\usepackage{hyperref}
\usepackage{xcolor}

\hypersetup{
  colorlinks=true,
  linkcolor=black,
  urlcolor=blue
}

\setlist[itemize]{leftmargin=2em,itemsep=0.25em,topsep=0.3em}
\setlist[enumerate]{leftmargin=2em,itemsep=0.25em,topsep=0.3em}
\linespread{1.18}

\newcommand{\slide}[2]{%
  \subsection*{第 #1 页：#2}
  \addcontentsline{toc}{subsection}{第 #1 页：#2}
}
\newcommand{\keyword}[1]{\textbf{\textcolor{blue!55!black}{#1}}}

\title{\bfseries 答辩逐字稿与问题速查}
\author{李思雨}
\date{2026年5月}

\begin{document}
\maketitle
\tableofcontents
\newpage

\section*{使用说明}
\addcontentsline{toc}{section}{使用说明}

这份材料用于答辩前演练和现场快速回忆。前半部分按 slides 顺序给出逐字稿，正式答辩时可以自然口语化，不必逐字背诵；后半部分整理了可能被问到的问题和建议回答。建议总时长控制在 12--15 分钟，理论页可以略快，模拟和实证页重点讲结论。

\newpage

\section{答辩逐字稿}

\slide{1}{标题页}
各位老师好，我是数学科学学院的李思雨。我的论文题目是《基于累积分布函数的非平稳空间序列统计推断》，指导老师是张园园老师。下面我将从研究动机、估计方法、主要理论结果、模拟研究、实际数据分析以及总结这几个方面进行汇报。

\slide{2}{目录}
本次汇报主要分为六个部分。首先介绍为什么要研究非平稳空间序列的误差分布，以及本文的总体思路；然后介绍模型、BST 趋势估计和 KDE 误差分布估计；接着说明本文的假设条件和主要理论结论；最后通过模拟和两个实际数据例子说明方法的表现，并给出总结和未来展望。

\slide{3}{研究动机与思路}
本文关注的是复杂空间区域上的非平稳空间序列。传统空间统计方法很多是建立在规则网格或矩形区域上的，但实际图像数据常常具有复杂边界，甚至会有空洞或不规则区域，因此直接套用规则区域方法会比较困难。

另外，误差分布推断在构造置信带、预测区间和模型诊断中都很重要。但已有许多误差分布方法假设观测独立，而图像或空间数据通常存在空间相关性。基于这两个问题，本文采用三角剖分上的二元样条来处理复杂区域，同时在空间相关条件下研究残差误差分布估计的 oracle 有效性，并进一步构造同时置信带。

\slide{4}{估计方法：模型}
本文考虑的基本模型是
\[
Y_{\mathbf{i}}=g(X_{\mathbf{i}})+Z_{\mathbf{i}}.
\]
其中 \(g\) 是未知的非平稳趋势函数，\(Z_{\mathbf{i}}\) 是严格平稳、零均值的空间误差过程。本文真正关心的是误差项的边际分布函数 \(F(z)\)。

为了减少趋势估计和误差分布估计之间的相互影响，本文采用样本分割的思想，把原始指标集分成两个互不相交的子集。一个子集用于构造残差并估计误差分布，另一个子集用于估计趋势函数。这样做的好处是理论推导更清晰，也能降低趋势估计误差对分布估计的干扰。

\slide{5}{估计方法：BST 数据划分与样条近似}
在复杂区域上，本文使用三角剖分上的二元样条，也就是 BST 方法。核心思想是把不规则区域划分成许多三角形，然后在每个三角形上用局部多项式来逼近趋势函数，并通过连续性约束保证整体光滑。

具体地，趋势函数 \(g(\boldsymbol{x})\) 可以用基函数展开近似为
\[
g(\boldsymbol{x})\approx \sum_{l\in\mathcal L}B_l(\boldsymbol{x})\gamma_l.
\]
这里 \(B_l\) 是样条基函数，\(\gamma_l\) 是待估参数。随后通过最小二乘方法在样条空间中估计 \(g\)。

\slide{6}{估计方法：BST 约束最小二乘}
由于样条函数需要满足跨三角形边界的连续性约束，所以原始最小二乘问题带有线性约束，可以写成
\[
\min_{\boldsymbol{\gamma}}\sum_{\mathbf{i}''}\{Y_{\mathbf{i}''}-\mathbf B^\top(\boldsymbol{x}_{\mathbf{i}''})\boldsymbol{\gamma}\}^2,
\quad \text{s.t. } \mathbf H\boldsymbol{\gamma}=0.
\]
这里的 \(\mathbf H\boldsymbol{\gamma}=0\) 表示样条系数需要满足的光滑约束。为了求解这个约束问题，本文对 \(\mathbf H^\top\) 做 QR 分解，把约束参数化为无约束参数，从而转化为普通最小二乘问题。

\slide{7}{估计方法：BST 无约束表示}
通过 QR 分解后，可以令 \(\boldsymbol{\gamma}=\mathbf Q_2\boldsymbol{\theta}\)，于是约束最小二乘转化为关于 \(\boldsymbol{\theta}\) 的无约束最小二乘。由此可以得到 \(\hat{\boldsymbol{\theta}}\) 和 \(\hat{\boldsymbol{\gamma}}\) 的显式表达。

最终，BST 趋势估计函数可以写为
\[
\hat g(\boldsymbol{x})=\mathbf B^\top(\boldsymbol{x})\hat{\boldsymbol{\gamma}}
=\mathbf P(\boldsymbol{x})\mathbf Y.
\]
这个表达在后续理论推导中很重要，因为它把趋势估计写成观测向量的线性组合。

\slide{8}{估计方法：趋势估计分解}
将 \(\hat g\) 的表达代入模型后，可以把趋势估计分解为两部分：
\[
\hat g(\boldsymbol{x})=\tilde g(\boldsymbol{x})+\tilde Z(\boldsymbol{x}).
\]
其中 \(\tilde g\) 是由真实趋势函数导致的近似部分，\(\tilde Z\) 是由空间误差传播到趋势估计中的随机部分。

这个分解是后面证明的关键。因为我们可以分别控制样条逼近误差和随机误差项，从而得到残差 \(\hat Z\) 与真实误差 \(Z\) 的一致接近性。

\slide{9}{估计方法：不可行 KDE 与 SCB}
如果真实误差 \(Z_{\mathbf{i}'}\) 可以观测，那么可以构造不可行的核分布估计量 \(\tilde F_{n',h}\)。它通过对真实误差进行核平滑，再积分得到分布函数估计。

在此基础上，可以构造光滑同时置信带：
\[
\left[\tilde F_{n',h}(z)-\frac{L_{1-\alpha}}{\sqrt{n'}},
\tilde F_{n',h}(z)+\frac{L_{1-\alpha}}{\sqrt{n'}}\right]\cap[0,1].
\]
但真实误差在实际中不可观测，所以这个估计量主要作为理论比较对象。

\slide{10}{估计方法：oracle KDE 与 SCB}
实际可用的是基于残差 \(\hat Z_{\mathbf{i}'}\) 的 oracle 核分布估计量 \(\hat F_{n',h}\)。这里所谓 oracle，不是说它知道真实误差，而是指它希望达到与不可行估计量相同的渐近效率。

本文的核心目标之一就是证明：用估计残差构造的 \(\hat F_{n',h}\)，在一致意义下与用真实误差构造的 \(\tilde F_{n',h}\) 渐近等价。因此它可以用于实际构造误差分布的光滑同时置信带。

\slide{11}{估计方法：非光滑 EMDF}
作为对照，本文还考虑了非光滑经验边际分布函数
\[
\hat F_{n'}(z)=n'^{-1}\sum_{\mathbf{i}'\in\mathcal I_{n'}}I(\hat Z_{\mathbf{i}'}\le z).
\]
它不使用核平滑，形式更直接。相应地，也可以把光滑同时置信带中的 \(\hat F_{n',h}\) 替换成 \(\hat F_{n'}\)，得到非光滑置信带。本文把它作为模拟研究中的基准方法，用来比较光滑方法的表现。

\slide{12}{理论假设：A1--A3}
本文的理论假设主要分为三类。第一类是关于误差分布光滑性的假设，即 \(F\) 和密度函数 \(f\) 需要有一定正则性。第二类是关于趋势函数 \(g\) 的光滑性假设，保证它可以被三角剖分样条有效逼近。

第三类是关于空间误差过程的假设。本文假设误差过程严格平稳、满足 \(\alpha\)-混合条件，并具有有限的 \(2+\delta\) 阶矩。这是处理空间相关性的关键条件，保证经验过程和核估计相关的概率界可以成立。

\slide{13}{理论假设：A4--A6}
后面三个假设分别对应三角剖分、核函数和带宽。A4 要求三角剖分是准均匀的，并规定三角形数量和样本量之间的增长关系。A5 要求核函数具有足够光滑性和紧支撑。A6 给出了带宽 \(h\) 的收敛速度范围。

这些假设的作用是平衡三类误差：样条逼近误差、趋势估计随机误差以及核平滑误差。只有它们在合适的速度下共同趋于零，后面的 oracle 有效性结论才能成立。

\slide{14}{主要理论结果：命题 1 和命题 2}
这一页给出理论推导中的几个中间结论。首先，推论说明不可行密度估计和核导数平均项在几乎处处意义下是有界的。这个结果用于后面对核函数差异项做 Taylor 展开和一致控制。

命题 1 控制的是样条逼近误差，即 \(\tilde g\) 与真实趋势 \(g\) 的差距，阶为 \(|\triangle|^{\ell+1}\)。命题 2 控制的是随机误差传播项 \(\tilde Z\)，阶为
\[
O_p\{n''^{-1/2}|\triangle|^{-1}(\log n)^{1/2}\}.
\]
这两个命题分别对应确定性逼近误差和随机估计误差。

\slide{15}{主要理论结果：命题 3}
由命题 1 和命题 2 可以推出命题 3。它说明真实误差 \(Z_{\mathbf{i}'}\) 与残差 \(\hat Z_{\mathbf{i}'}\) 的最大差异可以被 \(\|\hat g-g\|_\infty\) 控制，最终阶为
\[
O_p\{|\triangle|^{\ell+1}+n''^{-1/2}|\triangle|^{-1}(\log n)^{1/2}\}.
\]
直观来说，只要趋势函数估计得足够好，那么残差就能很好地逼近真实误差，这是用残差估计误差分布的基础。

\slide{16}{主要理论结果：定理一}
本文最核心的理论结果是定理一。在假设 A1 到 A6 成立时，基于残差的 oracle 估计量 \(\hat F_{n',h}\) 与基于真实误差的不可行估计量 \(\tilde F_{n',h}\) 渐近等价：
\[
D(\hat F_{n',h},\tilde F_{n',h})
=\|\hat F_{n',h}-\tilde F_{n',h}\|_\infty
=o_p(n'^{-1/2}).
\]
这说明趋势估计带来的残差误差不会影响一阶渐近性质，因此 \(\hat F_{n',h}\) 具有 oracle 有效性。

\slide{17}{置信区间：经验过程收敛}
为了构造同时置信带，需要进一步使用空间经验过程的收敛结论。在指数衰减混合条件和足够稀疏的采样下，经验分布函数会收敛到布朗桥。

结合定理一和已有关于平滑分布估计的结果，可以得到
\[
\|\hat F_{n',h}-F_{n'}\|_\infty=o_p(n'^{-1/2}).
\]
这一步说明用残差构造的光滑估计量在构造同时置信带时具有正确的渐近尺度。

\slide{18}{置信区间：SCB 渐近结论}
进一步可以得到
\[
\sqrt{n'}(\hat F_{n',h}-F)\Rightarrow \mathcal B(F).
\]
也就是说，残差光滑分布估计量的极限过程是布朗桥。由此可以构造 Kolmogorov--Smirnov 型的光滑同时置信带。

这里的 \(L_{1-\alpha}\) 是极限分布对应的分位数。实际理解上，就是用一个统一的误差半径同时覆盖所有 \(z\) 上的真实分布函数 \(F(z)\)。

\slide{19}{模拟研究：设计}
模拟研究中，本文分别考虑规则矩形区域和不规则 shape 区域。数据由模型
\[
Y_{\mathbf i}=g(x_1,x_2)+Z_{\mathbf i}
\]
生成，其中趋势函数为 \(\sin(\pi x_1)+\exp(-x_2^2)\)，误差项来自二维空间移动平均模型。

样本量设置为 \((100,50)\)、\((200,100)\)、\((200,200)\) 和 \((500,500)\)。模拟的主要目的有两个：一是比较不同估计量的最大偏差；二是考察同时置信带的实际覆盖率。

\slide{20--21}{模拟研究：三角剖分示例}
这两页展示了矩形区域和 shape 区域上的三角剖分示例。可以看到，三角剖分方法不仅适用于规则矩形，也可以自然适用于不规则边界。

这也是 BST 方法的优势：它不要求区域必须是矩形网格，而是通过局部三角形单元来刻画复杂区域，从而为后续趋势估计和残差分析提供统一框架。

\slide{22}{模拟研究：平均最大偏差}
这张表比较了三种估计量的平均最大偏差：光滑残差估计 \(\hat F_{n',h}\)、非光滑经验分布 \(\hat F_{n'}\)，以及不可行估计 \(\tilde F_{n',h}\)。

从结果看，随着样本量增大，三种估计量的偏差整体下降，说明估计具有一致性。并且 \(\hat F_{n',h}\) 与不可行估计 \(\tilde F_{n',h}\) 的表现非常接近，这从数值上支持了 oracle 有效性结论。

\slide{23--25}{模拟研究：覆盖率表}
这几页报告了不同区域和不同置信水平下同时置信带的覆盖率。总体来看，光滑方法 \(\hat F_{n',h}\) 的覆盖率接近名义水平，尤其在样本量增大后表现更稳定。

非光滑方法可以作为基准，但在一些情形下波动更明显。不可行估计 \(\tilde F_{n',h}\) 通常代表理论上更理想的参照，而 \(\hat F_{n',h}\) 与它接近，说明残差带来的额外误差在模拟中确实较小。

\slide{26}{模拟研究：箱线图}
箱线图展示的是统计量 \(D(\hat F_{n',h},F)\) 的分布情况。可以看到，随着样本量增加，箱线图整体向下收缩，离散程度也减小。

这说明估计误差不仅均值下降，而且稳定性增强。矩形区域和不规则区域都呈现类似趋势，进一步说明该方法对复杂区域具有适应性。

\slide{27}{模拟研究：99\% 同时置信带}
这一页直观展示了光滑 99\% 同时置信带与真实分布之间的比较。可以看到，在不同区域和不同样本量下，真实分布大体落在置信带内。

随着样本量增加，置信带明显变窄，这符合理论中半径为 \(n'^{-1/2}\) 量级的结论。也就是说，样本越多，对误差分布函数的不确定性刻画越精确。

\slide{28}{实际数据分析：数据来源}
实证部分使用两组真实数据。第一组是全球大气二氧化碳浓度数据，格点覆盖全球范围。第二组是黑海海表温度数据，区域具有明显的不规则边界。

选择这两组数据的原因是：CO\(_2\) 数据适合检验边际误差分布与常见参数分布之间的偏离；黑海数据则可以展示复杂不规则区域上方法的适用性。

\slide{29}{实际数据：全球 CO\texorpdfstring{$_2$}{2}}
这一页展示 CO\(_2\) 数据分析。左上图是 20 年平均浓度，右上图是趋势估计 \(\hat g\)。下方两幅图是基于误差分布的正态性检验结果。

主要结论是，传统正态模型在某些置信水平下可能无法充分描述实际误差分布。本文提出的同时置信带可以用来检验参数模型与真实边际误差分布之间的偏离程度。

\slide{30}{实际数据：黑海采样与剖分}
这一页展示黑海数据的空间采样位置、两个子样本的划分以及黑海区域的三角剖分。黑海区域并不是简单矩形，因此直接使用规则区域方法会比较不自然。

通过三角剖分，可以把复杂边界转化为若干三角形单元，并在其上进行 BST 趋势估计。这说明本文方法适用于真实复杂地理区域。

\slide{31}{实际数据：黑海海表温度}
这一页比较了 2020 年 12 月 9 日和 12 月 24 日的黑海海表温度，并展示了边际误差分布的 95\% 和 99\% 光滑同时置信带。

从图中可以看到，本文方法能够给出误差分布层面的不确定性评估。相比只看温度场本身，误差分布分析可以进一步帮助我们理解模型拟合后的残差结构。

\slide{32}{总结}
最后总结一下本文工作。本文针对非平稳空间序列，在指数衰减混合条件下构造了边际误差分布的光滑核分布估计量。

理论上，本文证明了基于残差的估计量在一致意义下具有 oracle 有效性，并建立了渐近 Kolmogorov--Smirnov 型同时置信带。数值模拟和两个实际数据分析说明，该方法可以用于复杂区域上的误差分布推断，并可用于评估传统参数模型的偏离程度。

未来可以进一步拓展到空间变系数模型，研究预测误差分布估计，并构建更精确的预测区间。

\slide{33--34}{参考文献}
这两页列出了本文主要参考文献，主要包括三角剖分样条、空间过程经验分布、误差分布估计和同时置信带等相关工作。它们分别支撑了本文的方法构造、理论证明和实际应用。

\slide{35}{致谢}
我的汇报到这里结束。感谢各位老师的聆听，恳请各位老师批评指正。

\newpage
\section{可能会被问到的问题速查}

\subsection*{1. 为什么要研究误差分布函数，而不是只估计均值函数？}
误差分布函数可以刻画模型无法解释的随机波动结构。只估计均值函数只能说明趋势层面的变化，而误差分布可以进一步用于模型诊断、预测区间构造和不确定性评估。本文的同时置信带本质上就是给误差分布函数提供统一的置信范围。

\subsection*{2. 为什么使用三角剖分上的二元样条？}
因为实际空间区域常常是不规则的，可能具有复杂边界或空洞。传统矩形网格方法对这种区域不够灵活。三角剖分可以自然贴合不规则边界，二元样条又能在每个三角形上进行局部多项式逼近，并通过连续性约束保证整体光滑。

\subsection*{3. 为什么要做样本分割？}
样本分割可以降低趋势估计和误差分布估计之间的依赖关系，使理论分析更清晰。一个子集用于估计趋势函数，另一个子集用于构造残差和估计误差分布。这样有助于证明残差分布估计与不可行真实误差分布估计的渐近等价。

\subsection*{4. 这里的 oracle 有效性是什么意思？}
oracle 有效性指的是：虽然真实误差 \(Z\) 不可观测，我们只能用残差 \(\hat Z\)，但基于残差构造的估计量 \(\hat F_{n',h}\) 在渐近意义下与基于真实误差的不可行估计量 \(\tilde F_{n',h}\) 等价。也就是说，使用估计残差不会损失一阶渐近效率。

\subsection*{5. 为什么需要核平滑？非光滑经验分布不可以吗？}
非光滑经验分布当然可以作为基准，但它是阶梯函数，不够平滑。核平滑后得到的分布估计更适合构造光滑同时置信带，也更便于理论推导和图形展示。模拟中也比较了光滑和非光滑方法，光滑方法整体表现稳定。

\subsection*{6. 带宽 \(h\) 如何选择？}
理论上，带宽需要满足假设 A6 中的收敛速度，以平衡核平滑偏差和随机误差。实际模拟中采用适合样本规模的带宽规则。答辩时可以强调：本文重点是建立在满足理论带宽条件下的渐近性质，实际带宽选择可以作为后续进一步研究的问题。

\subsection*{7. 为什么要求误差过程满足 \(\alpha\)-混合？}
\(\alpha\)-混合条件用于刻画空间相关性随距离增加而衰减。它比独立假设更适合空间数据，同时又足够强，可以保证经验过程收敛、核估计一致性和相关概率不等式成立。

\subsection*{8. 定理一证明的核心思路是什么？}
核心是控制
\[
\|\hat F_{n',h}-\tilde F_{n',h}\|_\infty.
\]
通过 Taylor 展开把核函数差异转化为残差误差 \(\hat Z-Z\) 的影响，再用命题 3 控制残差和真实误差的最大距离，同时利用核函数导数项的有界性，最终得到该差异是 \(o_p(n'^{-1/2})\) 量级。

\subsection*{9. 命题 1、2、3 之间是什么关系？}
命题 1 控制确定性的样条逼近误差，命题 2 控制趋势估计中的随机误差传播。两者合起来给出 \(\|\hat g-g\|_\infty\) 的阶。命题 3 再利用
\[
|Z_{\mathbf{i}'}-\hat Z_{\mathbf{i}'}|\le \|\hat g-g\|_\infty
\]
得到残差与真实误差的最大差异界。

\subsection*{10. 同时置信带里的 \(L_{1-\alpha}\) 是什么？}
\(L_{1-\alpha}\) 是极限 Kolmogorov--Smirnov 型分布的分位数，可以理解为布朗桥上确界统计量的临界值。它使得整个函数 \(F(z)\) 在所有 \(z\) 上同时被覆盖，而不是只在某个固定点上覆盖。

\subsection*{11. 为什么置信带要和 \([0,1]\) 取交集？}
因为分布函数的取值必须在 \([0,1]\) 内。直接加减置信半径后，端点附近可能小于 0 或大于 1，所以需要截断到合法的分布函数范围内。

\subsection*{12. 模拟中三种估计量分别代表什么？}
\(\hat F_{n',h}\) 是本文实际可用的光滑残差估计量；\(\hat F_{n'}\) 是非光滑经验分布函数，用作基准；\(\tilde F_{n',h}\) 是基于真实误差的不可行估计量，用作理论参照。比较三者可以检验 residual-based 方法是否接近 oracle 情形。

\subsection*{13. 覆盖率表中星号是什么意思？}
星号用于标记不同方法之间覆盖率比较的结果。答辩时如果老师追问，可以说明：表格主要用来展示覆盖率是否接近名义水平，以及光滑方法、非光滑方法和不可行方法之间的相对表现；星号是辅助比较符号，不影响主要理论结论。

\subsection*{14. 实证数据为什么选 CO\(_2\) 和黑海 SST？}
CO\(_2\) 数据覆盖全球规则格点，适合展示误差分布与正态模型之间的偏离检验。黑海海表温度数据具有复杂地理边界，适合展示三角剖分和 BST 方法在不规则区域上的适用性。两者分别体现了分布诊断和复杂区域建模两个侧重点。

\subsection*{15. 本文方法的局限性是什么？}
主要有三点。第一，理论依赖较强的混合条件和光滑性条件；第二，带宽、三角剖分密度和样本分割比例在实际中需要调节；第三，目前主要研究边际误差分布，尚未进一步处理条件误差分布或更复杂的时空动态结构。

\subsection*{16. 未来工作可以怎么做？}
可以拓展到空间变系数模型，研究预测误差分布估计和预测区间构造。也可以进一步研究自适应带宽选择、非均匀采样下的理论性质，以及更复杂的时空相关结构。

\subsection*{17. 如果老师问“你的创新点是什么”？}
可以回答：本文的创新点主要有三点。第一，把复杂区域上的 BST 趋势估计与误差分布函数推断结合起来；第二，在空间相关条件下证明了基于残差的核分布估计量具有 oracle 有效性；第三，构造了适用于复杂区域非平稳空间序列的光滑同时置信带，并通过模拟和真实数据验证了方法效果。

\subsection*{18. 如果老师问“论文中最关键的技术难点是什么”？}
可以回答：最关键的是控制趋势估计误差对误差分布估计的影响。由于误差不可观测，我们只能用残差，而残差中包含趋势估计误差。因此需要先证明 \(\hat g\) 与 \(g\) 的一致接近，再证明这种接近性不会影响核分布估计的一阶渐近性质。

\subsection*{19. 如果老师问“样本分割会不会浪费数据”？}
可以回答：样本分割确实会牺牲一部分样本利用效率，但它换来了理论上的独立性或弱依赖结构简化，使 oracle 有效性证明更可行。在样本量较大时，这种损失可以通过合理设置 \(n'\)、\(n''\) 和 \(T_n\) 来控制。未来也可以研究交叉拟合来提高数据利用率。

\subsection*{20. 如果老师问“为什么不是直接 bootstrap”？}
可以回答：空间数据存在相关性，普通 bootstrap 不一定适用。空间 block bootstrap 或 subsampling 是可选方法，但理论和实现会更复杂。本文选择基于经验过程极限和核平滑的解析式置信带，重点是建立可解释的渐近理论。

\newpage
\section{一分钟压缩版总结}

本文研究复杂区域上非平稳空间序列的误差分布推断。方法上，先用三角剖分上的二元样条估计非平稳趋势，再用残差构造核分布估计量和光滑同时置信带。理论上，证明了残差估计量与基于真实误差的不可行估计量渐近等价，即具有 oracle 有效性，并进一步得到 Kolmogorov--Smirnov 型同时置信带。模拟结果表明估计误差随样本量增大而下降，覆盖率接近名义水平；CO\(_2\) 和黑海海表温度数据说明该方法可用于误差分布诊断和复杂区域上的实际空间数据分析。

\end{document}
